\documentclass[10pt,a4paper]{amsart}
\usepackage[margin=19mm]{geometry}
\usepackage{amsmath,amssymb,mathtools}
\usepackage[scr=rsfs,cal=euler]{mathalfa}
\usepackage{xfrac}
\usepackage[unicode,hidelinks,bookmarks=false]{hyperref}
\newcommand{\ie}{i.e.\ }
\newcommand{\cf}{cf.\ }
\newcommand{\resp}{respectively\ }
\newcommand{\Subsection}{\subsection}
\providecommand{\nobreakdash}{\nobreak\hbox{-}}
\setlength{\emergencystretch}{2em}
\multlinegap0pt
\usepackage{bm}
\usepackage{bbm}
\usepackage{dsfont}
\usepackage{xspace}
\usepackage{cancel}
\usepackage{mathtools}
\usepackage{amsthm}
\makeatletter
\@ifundefined{thm}{\newtheorem{thm}{Thm}}{}
\@ifundefined{cor}{\newtheorem{cor}[thm]{Corollary}}{}
\@ifundefined{defn}{\newtheorem{defn}{Definition}}{}
\@ifundefined{lemma}{\newtheorem{lemma}[thm]{Lemma}}{}
\@ifundefined{prop}{\newtheorem{prop}[thm]{Proposition}}{}
\makeatother
\newcommand{\Z}{\mathbb{Z}}
\newcommand{\dc}{\,\underline{\mathrm{d}}^c}
\newcommand{\mc}{\, \underline{\mathrm{m}}^c}
\newcommand{\nr}{\mathrm{nr}}
\newcommand{\nw}{\mathrm{nw}}
\newcommand{\w}{\mathrm{w}}
\title[Classification of marked elliptic root systems with non-redu]{Classification of marked elliptic root systems with non-reduced quotient}
\author{A. Fialowski, K. Iohara, Y. Saito}
\date{Source: arXiv:2408.01358v2 (CC BY 4.0)}
\begin{document}
\maketitle

\noindent\emph{Source and licence.} Classification of marked elliptic root systems with non-reduced quotient. Version arXiv:2408.01358v2 (CC BY 4.0), distributed under the Creative Commons Attribution 4.0 International licence, as recorded in the arXiv licence metadata for this exact version and shown on the abstract page https://arxiv.org/abs/2408.01358v2. Full rights notice: licenses/arxiv-2408.01358-CC-BY-4.0.txt.\par\smallskip

\noindent\textit{Editorial scope.} Counted unit: the Lemma whose source label is \texttt{lemma:Gamma(R'',G)} in the article (source0.tex lines 9476--9499, source label \texttt{lemma:Gamma(R'',G)}), delivered together with the article's own complete original proof of it (source0.tex lines 9501--9663, one \texttt{proof} environment of 163 lines). No mathematics is written, repaired or re-derived: statement and proof are the article's own text line for line. Required context retained uncounted: k=kast (lines 1165--1167); defn:nr-countings\_alpha'-ast (lines 6086--6107); cor:nr-countings1 (lines 6322--6339); prop:countings-prime (lines 6392--6407); lemma:black-ast (lines 9101--9109); eq:description-Rmc (lines 9394--9398); eq:defn-(alpha)R'' (lines 9433--9438); eq:a-Rmc (lines 9468--9470). Every definition, display and label of those lines is retained unchanged. Omitted from this package and not redistributed: 1--1164, 1168--6085, 6108--6321, 6340--6391, 6408--9100, 9110--9393, 9399--9432, 9439--9467, 9471--9475, 9500--9500, 9664--18375. Nothing inside a retained excerpt is omitted or altered: every retained body is delivered exactly as the article's lines are written, so no omission receipt of any kind accompanies this package. Citations: the retained excerpts carry no citation command at all, so no bibliography entry of the article is transcribed and no reference is redistributed. Reference adaptation: none was needed and none was made; every reference inside the delivered unit resolves inside it, so no unresolved reference or citation is produced. Printed slips kept literal: no printed word, symbol, hypothesis or number of the counted statement or of its proof was altered. Layout and class: the article's own class is \texttt{amsart} with packages including latexsym, graphicx, amssymb, xy, etex, tikz, diagbox, adjustbox, fontenc, xcolor, ulem, hyperref, setspace, imakeidx; the wrapper is the lane's pinned \texttt{amsart} editorial wrapper at 10pt on A4, which restates the theorem-like environments the retained text needs and adds the article's own macro definitions that the retained text needs (\texttt{\textbackslash Z}, \texttt{\textbackslash dc}, \texttt{\textbackslash mc}, \texttt{\textbackslash nr}, \texttt{\textbackslash nw}, \texttt{\textbackslash w}), with extra wrapper packages (bm, bbm, dsfont, xspace, cancel, mathtools). The delivered numbering is the unit's own. Assets: no retained span contains a figure environment, a graphics-inclusion command, a PDF-page inclusion, a table or a TikZ picture; no image file is copied into this package, the package declares no asset dependency and no image file is redistributed with it. Licence: version arXiv:2408.01358v2 of this article is distributed under the Creative Commons Attribution 4.0 International licence, as recorded in the arXiv licence metadata for this exact version and shown on the abstract page https://arxiv.org/abs/2408.01358v2. A full rights notice is delivered at licenses/arxiv-2408.01358-CC-BY-4.0.txt. Characters: every retained excerpt is pure ASCII, so the delivered engine, XeLaTeX with its default Latin Modern text font, reports no missing character.
\begin{equation}\label{k=kast}
k(\alpha)=k(\alpha^\ast).
\end{equation}

\begin{defn}
For an element $\alpha \in \Pi^{\dc}$, define a positive integer
$k^{\nr}((\alpha')^\ast)\in \Z_{>0}$\index[notations]{k711@$k^{\nr}((\alpha')^\ast)$} and 
a positive half integer 
$k^{\nr}(\alpha)\in\frac12 \Z_{>0}$ by the next equalities{\rm :}
\begin{align} 
(\alpha')^\ast=&((\alpha')^\ast:\alpha)_G\,\alpha + k^{\nr}((\alpha')^\ast)a, 
\label{defn:nr-countings_alpha'-ast}
\\
\alpha=&(\alpha:(\alpha')^\ast)_G\,(\alpha')^\ast-k^{\nr}(\alpha)a,
\label{defn:nr-countings_alpha}
\end{align}
where 
\begin{equation*}
(x:y)_G:=(\overline{x}:\overline{y}).
\end{equation*}
We call $k^{\nr}(\alpha)$\index[notations]{k712@$k^{\nr}(\alpha)$} the 
\textbf{non-reduced counting number}\index[index]{counting number!non-reduced 
counting number@non-reduced --}  
%  for $\rightarrow$ 
of an element $\alpha\in \Pi^{\dc}\cup (\Pi^{\mc})^\ast$.
\end{defn}

\begin{cor}\label{cor:nr-countings1}
For $\alpha\in \Pi^{\dc}$, we have
\begin{align}
k^{\mathrm{nr}}((\alpha')^\ast)&=
\begin{cases}
\frac12 k((\alpha')^\ast) & \text{if }2\alpha\not\in R\text{ and }
2\overline{\alpha}\in R/G,\\
k((\alpha')^\ast) & \text{otherwise},
\end{cases}\label{nr-countings_alpha3}  
\\
k^{\mathrm{nr}}(\alpha)&=
\begin{cases}
\frac12 k(\alpha) & \text{if }\frac12(\alpha')^\ast\not\in R\text{ and }
\frac12\overline{(\alpha')^\ast}\in R/G,\\
k(\alpha) & \text{otherwise}.
\end{cases}\label{nr-countings_alpha4}
\end{align}
\end{cor}

\begin{prop}\label{prop:countings-prime}
Let $\alpha\in \Pi^{\dc}$. If $2\overline{\alpha}\in R/G\ \left(\Longleftrightarrow\ 
\frac12 \overline{(\alpha')^\ast}\in R/G\right)$, we have 
\begin{equation}\label{countings-alpha-prime}
k((\alpha')^\ast)=\begin{cases}
2k(\alpha) & \text{if } 2\alpha\not\in R\text{ and } \frac12 (\alpha')^\ast\not\in R,\\
4k(\alpha) & \text{if } 2\alpha\not\in R\text{ and }\frac12 (\alpha')^\ast\in R,\\
k(\alpha) & \text{if } 2\alpha\in R\text{ and }\frac12 (\alpha')^\ast\not\in R,\\
2k(\alpha) & \text{if } 2\alpha\in R\text{ and }\frac12 (\alpha')^\ast\in R.
\end{cases}
\end{equation}
Otherwise, 
\begin{equation}
k((\alpha')^\ast)=k(\alpha). 
\end{equation}
\end{prop}

\begin{lemma}\label{lemma:black-ast}
For $\alpha_i\in |\Gamma^\downarrow|$, the following two conditions are equivalent.
\begin{itemize}
\item[(a)] $\frac12 (\alpha_i')^\ast\in R$.
\vskip 1mm
\item[(b)] $\alpha_i\in |\Gamma^\nw_m|_g^b\amalg |\Gamma^\nw_m|_b^b\amalg 
|(\Gamma^{\downarrow}\setminus \Gamma^{\downarrow}_m)^\nw|$.
\end{itemize}
\end{lemma}

\begin{equation}\label{eq:description-Rmc}
R^{\mc}=\Big(\bigcup_{\alpha_i\in |\Gamma^\w|}R^{\mc}(i)\Big)\bigcup
\Big(\bigcup_{\alpha_j\in |\Gamma^\nw|_b}R^{\mc}(j)\Big)\bigcup
\Big(\bigcup_{\alpha_j\in |\Gamma^\nw|_g}R^{\mc}(j)\Big).
\end{equation}

\begin{equation}\label{eq:defn-(alpha)R''}
(\alpha_i)_{R^{\mc}}=\begin{cases}
2\alpha_i & \text{if $\alpha_i\in |\Gamma^\nw|_b$},\\
\alpha_i & \text{otherwise}. 
\end{cases}
\end{equation}

\begin{equation}\label{eq:a-Rmc}
a^{R^{\mc}}=[Q(R)\cap G:Q(R^{\mc})\cap G]a.
\end{equation}

\begin{lemma}\label{lemma:Gamma(R'',G)}
{\rm (1)} Let $k_{R^{\mc}}(\alpha)$ be the counting number of 
$\alpha\in R^{\mc}$ as an element of the reduced mERS $(R^{\mc},G)$. One has
\begin{equation}\label{eq:counting-Rmc-1}
[Q(R)\cap G:Q(R^{\mc})\cap G]\,
k_{R^{\mc}}\big((\alpha_i)_{R^{\mc}}\big)=\begin{cases}
2k(\alpha_i) & \text{if $\frac12 (\alpha_i')^\ast\in R$},\\
k(\alpha_i) & \text{otherwise},
\end{cases}
\end{equation}
\begin{equation}\label{eq:counting-Rmc-2}
[Q(R)\cap G:Q(R^{\mc})\cap G]\,
k_{R^{\mc}}\big(\big((\alpha_i)_{R^{\mc}}'\big)^\ast\big)=
k(\alpha_i')\quad\text{for every}\quad 0\leq i\leq l, 
\end{equation}
{\rm (2)} The explicit forms of $\big((\alpha_i)_{R^{\mc}}'\big)^\ast\in 
\big(\Pi^{\mc}_{R^{\mc}}\big)^\ast\, (0\leq i\leq l)$ are given by
\[\big((\alpha_i)_{R^{\mc}}'\big)^\ast=
(\alpha_i')^\ast.\]
{\rm (3)} The explicit forms of 
$(\alpha_i)_{R^{\mc}}'\in \Pi^{\mc}_{R^{\mc}}\, (0\leq i\leq l)$ 
are given by
\[(\alpha_i)_{R^{\mc}}'=\alpha_i'.\]
\end{lemma}

\begin{proof}
(1) The formula \eqref{eq:counting-Rmc-1} is an easy consequence of the description 
\eqref{eq:description-Rmc} of $R^{\mc}$. Let us prove \eqref{eq:counting-Rmc-2}. 
Note that, since $(R^{\mc},G)$ is reduced, there is no black node in 
$|\Gamma(R^{\mc},G)|$. 
Applying Proposition \ref{prop:countings-prime} for $(R^{\mc},G)$, 
we have
\begin{equation}\label{eq:couting-Rmc-ast}
\begin{split}
&k_{R^{\mc}}\big(\big((\alpha_i)_{R^{\mc}}'\big)^\ast\big) \\
=
&
\begin{cases}
k_{R^{\mc}}\big((\alpha_i)_{R^{\mc}}\big) & 
\text{if $(\alpha_i)_{R^{\mc}}$ is a white node in $|\Gamma(R^{\mc},G)|$},\\
2k_{R^{\mc}}\big((\alpha_i)_{R^{\mc}}\big) & 
\text{if $(\alpha_i)_{R^{\mc}}$ is a gray node in $|\Gamma(R^{\mc},G)|$}.
\end{cases}
\end{split}
\end{equation}
By the definition of $(\alpha_i)_{R^{\mc}}$, we have
\begin{equation}\label{eq:color-Rmc}
\left\{\begin{array}{lllllll}
\text{$(\alpha_i)_{R^{\mc}}$ is a white node in $\Gamma(R^{\mc},G)$} 
& \Longleftrightarrow & \alpha_i\in |\Gamma^\w|\amalg |\Gamma^\nw|_b,\\[3pt]
\text{$(\alpha_i)_{R^{\mc}}$ is a gray node in $\Gamma(R^{\mc},G)$} 
& \Longleftrightarrow & \alpha_i\in |\Gamma^\nw|_g.
\end{array}\right.
\end{equation}
Therefore, by \eqref{eq:couting-Rmc-ast} and
\eqref{eq:color-Rmc}, we have
\begin{align*}
&[Q(R)\cap G:Q(R^{\mc})\cap G]\,
k_{R^{\mc}}\big(\big((\alpha_i)_{R^{\mc}}'\big)^\ast\big)\\
&\quad =[Q(R)\cap G:Q(R^{\mc})\cap G]\,
k_{R^{\mc}}\big((\alpha_i)_{R^{\mc}}\big)\times
\begin{cases}
1 & \text{if $\alpha_i\in |\Gamma^\w|\amalg |\Gamma^\nw|_b$},\\
2 & \text{if $\alpha_i\in |\Gamma^\nw|_g$}.
\end{cases}
\end{align*}
Furthermore, \eqref{eq:counting-Rmc-1} and Lemma \ref{lemma:black-ast}
imply that
\begin{align*}
&\text{the right hand side} =\begin{cases}
k(\alpha_i) & \text{if } \alpha_i\in |\Gamma^\w|\amalg |\Gamma_m^\nw|_b^g ,\\
2k(\alpha_i) & \text{if } \alpha_i\in |\Gamma_m^\nw|_g^g\amalg |\Gamma_m^\nw|_b^b
\amalg |(\Gamma^{\downarrow}\setminus\Gamma^{\downarrow}_m)^\nw|,\\
4k(\alpha_i) & \text{if }\alpha_i\in|\Gamma_m^\nw|_g^b.
\end{cases}
\end{align*}
By Proposition \ref{prop:countings-prime} for $(R,G)$, this coincides with 
$k((\alpha_i')^\ast)=k(\alpha_i')$. Thus, we have \eqref{eq:counting-Rmc-2}.
\vskip 1mm
\noindent
(2) Note that the formula
\begin{equation}\label{eq:n-red-coungting_in_Rmc}
k_{R^{\mc}}^{nr}\big(\big((\alpha_i)_{R^{\mc}}'\big)^\ast\big)=
k_{R^{\mc}}\big((\alpha_i)_{R^{\mc}}\big)
\end{equation}
holds. Indeed, by Corollary \ref{cor:nr-countings1} for $(R^{\mc},G)$, we have
\begin{align*}
&k_{R^{\mc}}\big(\big((\alpha_i)_{R^{\mc}}'\big)^\ast\big) \\
=
&\begin{cases}
k_{R^{\mc}}^{nr}\big(\big((\alpha_i)_{R^{\mc}}'\big)^\ast\big) & 
\text{if $(\alpha_i)_{R^{\mc}}$ is a white node in $\Gamma(R^{\mc},G)$},\\
2k_{R^{\mc}}^{nr}\big(\big((\alpha_i)_{R^{\mc}}'\big)^\ast\big) &
\text{if $(\alpha_i)_{R^{\mc}}$ is a gray node in $\Gamma(R^{\mc},G)$}.
\end{cases}
\end{align*}
This formula and \eqref{eq:couting-Rmc-ast} imply
\eqref{eq:n-red-coungting_in_Rmc}.

Hence, by \eqref{defn:nr-countings_alpha'-ast} and \eqref{eq:n-red-coungting_in_Rmc},
we get 
\begin{equation}\label{eq:prime-ast-Rmc}
\begin{aligned}
\big((\alpha_i)_{R^{\mc}}'\big)^*=
%\big(\big((\alpha_i)_{R^{\mc}}'\big)^*:(\alpha_i)_{R^{mc}}\big)_G(\alpha_i)_{R^{\mc}}
%+k_{R^{\mc}}^{nr}\big(\big((\alpha_i)_{R^{\mc}}'\big)^\ast\big)a^{R^{\mc}} \quad 
%(\because \ \eqref{defn:nr-countings_alpha'-ast})\\
\big(\big((\alpha_i)_{R^{\mc}}'\big)^*:(\alpha_i)_{R^{mc}}\big)_G(\alpha_i)_{R^{\mc}}
+k_{R^{\mc}}\big((\alpha_i)_{R^{\mc}}\big)a^{R^{\mc}}.
%(\because\ \eqref{eq:n-red-coungting_in_Rmc}).
\end{aligned}
\end{equation}
By \eqref{eq:color-Rmc} and the definition \eqref{eq:defn-(alpha)R''} 
of $(\alpha_i)_{R^{\mc}}$, we have
\begin{align*}
\big(\big((\alpha_i)_{R^{\mc}}'\big)^*:(\alpha_i)_{R^{mc}}\big)_G(\alpha_i)_{R^{\mc}}
&=\begin{cases}
(\alpha_i)_{R^{\mc}} & \text{if $(\alpha_i)_{R^{\mc}}$ is a white node},\\
2(\alpha_i)_{R^{\mc}} & \text{if $(\alpha_i)_{R^{\mc}}$ is a gray node},
\end{cases}\\
&=\begin{cases}
\alpha_i & \text{if $\alpha_i\in |\Gamma^\w|$},\\
2\alpha_i & \text{otherwise},
\quad 
\end{cases}
\end{align*}
Furthermore, by \eqref{eq:a-Rmc} and \eqref{eq:counting-Rmc-1}, we get 
%{\color{red} the next conditions sound strange...}
\begin{align*}
k_{R^{\mc}}\big((\alpha_i)_{R^{\mc}}\big)a^{R^{\mc}} 
&=[Q(R)\cap G:Q(R^{\mc})\cap G]k_{R^{\mc}}\big((\alpha_i)_{R^{\mc}}\big)a\\
%\quad (\because\ \eqref{eq:a-Rmc})\\
&=\begin{cases}
2k(\alpha_i)a & \text{if $\alpha_i\in 
|\Gamma_m^\nw|_g^b
\amalg |\Gamma_m^\nw|_b^b\amalg 
|(\Gamma^{\downarrow}\setminus 
\Gamma_m^{\downarrow})^\nw|$},\\
k(\alpha_i)a & \text{if $|\Gamma^\w|\amalg 
|\Gamma_m^\nw|_g^g\amalg |\Gamma_m^\nw|_b^g$}.
\end{cases}
\end{align*}
Substituting the above two formulas to \eqref{eq:prime-ast-Rmc}, we have
\begin{align*}
\big((\alpha_i)_{R^{\mc}}'\big)^*
&=
%&\qquad =\begin{cases}
%(\alpha_i)_{R^{\mc}}+k_{R^{\mc}}\big((\alpha_i)_{R^{\mc}}\big)a^{R^{\mc}} & 
%\text{if $(\alpha_i)_{R^{\mc}}$ is a white node in $\Gamma(R^{\mc},G)$},\\
%2(\alpha_i)_{R^{\mc}}+k_{R^{\mc}}\big((\alpha_i)_{R^{\mc}}\big)a^{R^{\mc}} & 
%\text{if $(\alpha_i)_{R^{\mc}}$ is a gray node in $\Gamma(R^{\mc},G)$},
%\end{cases}\\
%&\qquad =\begin{cases}
%(\alpha_i)_{R^{\mc}}+k(\alpha_i)a & \text{if $\alpha_i\in 
%|\Gamma^\w|\amalg |\Gamma_m^\nw|_b^g$},\\
%(\alpha_i)_{R^{\mc}}+2k(\alpha_i)a & \text{if $\alpha_i\in 
%|\Gamma_m^\nw|_b^b\amalg |(\Gamma^{\downarrow}\setminus 
%\Gamma_m^{\downarrow})^\nw|$},\\
%2(\alpha_i)_{R^{\mc}}+k(\alpha_i)a & \text{if $\alpha_i\in |\Gamma_m^\nw|_g^g$},\\
%2(\alpha_i)_{R^{\mc}}+2k(\alpha_i)a & \text{if $\alpha_i\in |\Gamma_m^\nw|_g^b$},
%\end{cases}\\
%&\qquad =
\begin{cases}
\alpha_i+k(\alpha_i)a & \text{if $\alpha_i\in |\Gamma^\w|$},\\
%\alpha_i+2k(\alpha_i)a & \text{if $\alpha_i\in |\Gamma_m^\nw|_b^g$},\\
2\alpha_i+k(\alpha_i)a & \text{if $\alpha_i\in 
|\Gamma_m^\nw|_b^g\amalg |\Gamma_m^\nw|_g^g$},\\
2\alpha_i+2k(\alpha_i)a & \text{if $\alpha_i\in 
|\Gamma_m^\nw|_g^b\amalg |\Gamma_m^\nw|_b^b \amalg
|(\Gamma^{\downarrow}\setminus \Gamma_m^{\downarrow})^\nw|$}.\\
\end{cases}
\end{align*}
This coincides with $(\alpha_i')^\ast$, as desired.
\vskip 1mm
\noindent
(3) 
By the definition of the $\ast$-operation and \eqref{k=kast}, we have
\[\alpha_i'=(\alpha_i')^\ast-k((\alpha_i')^\ast)a\quad{and}\quad
(\alpha_i)_{R^{\mc}}'
=\big((\alpha_i)_{R^{\mc}}'\big)^\ast
-k_{R^{\mc}}\big(\big((\alpha_i)_{R^{\mc}}'\big)^\ast\big)a^{R^{\mc}}.\]
Moreover, by statement (2), \eqref{eq:a-Rmc} and 
\eqref{eq:counting-Rmc-2}, we get 
\[\big((\alpha_i)_{R^{\mc}}'\big)^\ast=(\alpha_i')^\ast\quad\text{and}\quad
k_{R^{\mc}}\big(\big((\alpha_i)_{R^{\mc}}'\big)^\ast\big)a^{R^{\mc}}
=k((\alpha_i')^\ast)a.\]
Compiling these four formulas, we have the desired result.
 \end{proof}

\end{document}
